\documentclass[10pt,a4paper]{amsart}
\usepackage[margin=19mm]{geometry}
\usepackage{amsmath,amssymb,mathtools}
\usepackage[scr=rsfs,cal=euler]{mathalfa}
\usepackage{xfrac}
\usepackage[unicode,hidelinks,bookmarks=false]{hyperref}
\newcommand{\ie}{i.e.\ }
\newcommand{\cf}{cf.\ }
\newcommand{\resp}{respectively\ }
\newcommand{\Subsection}{\subsection}
\providecommand{\nobreakdash}{\nobreak\hbox{-}}
\setlength{\emergencystretch}{2em}
\multlinegap0pt
\usepackage{mathtools}
\usepackage{amssymb}
\usepackage{dsfont}
\usepackage{enumerate}
\usepackage{csquotes}
\usepackage{xcolor}
\usepackage{bm}
\newtheorem{theorem}{Theorem}
\newtheorem{lemma}{Lemma}
\def \Ghat{\widehat{G}}
\def\N{{\mathbb N}}% nonnegative integers
\def \O{\mathscr{O}}
\def\R{{\mathbb R}}% real numbers
\def\Tend#1#2{\mathop{\longrightarrow}\limits_{#1\rightarrow#2}}
\def\eps{\varepsilon}
\def \princ{{\rm princ}_{\hbar}}
\title{Quantum ergodicity for contact metric structures}
\author{Lino Benedetto}
\date{Source: arXiv:2507.18216v4 (CC BY 4.0)}
\begin{document}
\maketitle

\noindent\emph{Source and licence.} Quantum ergodicity for contact metric structures. Version arXiv:2507.18216v4 (CC BY 4.0), distributed by arXiv under the Creative Commons Attribution 4.0 International licence, http://creativecommons.org/licenses/by/4.0/, as recorded in the arXiv licence metadata for this exact version and shown on the abstract page https://arxiv.org/abs/2507.18216v4. Full licence notice: \texttt{licenses/arxiv-2507.18216-CC-BY-4.0.txt}.\par\smallskip

\noindent\textit{Editorial scope.} Counted unit: the article's theorem thm:microlocalweyllaw (source0.tex lines 2332--2339). The article's complete original proof of it is delivered with it (source0.tex lines 2396--2426, one \texttt{proof} environment of 31 lines, which the article places away from the statement). No mathematics is written, repaired or re-derived: the statement and the proof are the article's own lines, in the article's own notation, and no retained line is substituted or re-flowed.  Retained uncounted as the source-local context the proof needs: lines 1799--1810; lines 2351--2357.  Nothing was omitted from any retained span, so the package carries no omission record.  No retained line contains a citation, so the wrapper carries no bibliography and no bibliography entry.  No retained line contains a figure environment, an image-inclusion command or drawing code, so no image is delivered and the package declares no asset dependency.  Editorial formatting only, in addition: an AMS \texttt{amsart} preamble that restates the article's own declarations and macros used by the retained lines, namely the environments \texttt{theorem}, \texttt{lemma} on separate counters, together with the standard packages those lines need.  The retained environments print in this document as Theorem 1, Lemma 1 in order of appearance, whereas the article prints the same items with the numbering of its own full document.  The complete native source of the article as distributed in the arXiv e-print for this exact version is bundled unchanged as \texttt{sources/182136/source0.tex}.
\begin{theorem}
    \label{thm:functionalcalc}
    Let $f\in C_c^\infty(\R)$. For any $N\in \N$, there exist $Q_{f,N}(\hbar)\in \Psi_\hbar^{-\infty}(M)$ and $R_{f,N}(\hbar)$ a bounded operator on $L^2(M)$ such that
    \begin{equation*}
        f(-\hbar^2\Delta_{sR}) = Q_{f,N}(\hbar) + R_{f,N}(\hbar),\quad\mbox{with}\ R_{f,N}(\hbar) = \O_{L^2(M)\rightarrow L^2(M)}(\hbar^{N+1}).
    \end{equation*}
    Moreover, the principal symbol of $Q_{f,N}(\hbar)$ is given by
    \begin{equation*}
        \princ(Q_{f,N}(\hbar)) = f(H),
    \end{equation*}
    where $H$ denotes the principal symbol $\princ(-\hbar^2\Delta_{sR})$ of the subLaplacian.
\end{theorem}

\begin{lemma}
    \label{lem:traceclass}
    Any operator $B\in \Psi_\hbar^{-\infty}(M)$ is of trace class on $L^2(M)$. Moreover, if $\beta\in S^{-\infty}(\Ghat M)$ denotes the principal symbol of $B$, then we have
    \begin{equation*}
        {\rm Tr}(B) = {\rm vol}(M)\int_{\Ghat M} {\rm Tr}_{\mathcal{H}_\pi}(\beta(x,\pi))\,d\hat{\mu}_{\Ghat M}(x,\pi) + \O(\hbar).
    \end{equation*}
\end{lemma}

\begin{theorem}
    \label{thm:microlocalweyllaw}
    Let $B\in \Psi_\hbar^{-\infty}(M)$ be an admissible semiclassical pseudodifferential operator. Then, for any two real numbers $a < b$, we have the following asymptotic
    \begin{equation*}
        \hbar^{Q} \sum_{a\leq \delta_k(\hbar)\leq b}\langle B\varphi_k,\varphi_k\rangle \Tend{\hbar}{0} {\rm vol}(M)\int_{\Ghat M} {\rm Tr}\left({\bm 1}_{[a,b]}(H)(x,\pi)\beta(x,\pi)\right) d\mu_{\Ghat M}(x,\pi),
    \end{equation*}
    where $\beta \in S^{-\infty}(\Ghat M)$ denotes the principal symbol of $B$ and we recall that $Q = 2d+2$ denotes the \emph{homogeneous dimension} of $(M,\eta)$.
\end{theorem}

\begin{proof}[Proof of Theorem~\ref{thm:microlocalweyllaw}]
    Let $\eps>0$ and let $\underline{f_\eps}$ and $\overline{f_\eps}$ be two smooth functions on $\R$ valued in $[0,1]$ satisfying
    \begin{equation*}
        \underline{f} {\bm 1}_{[a,b]} = \underline{f}\quad \mbox{and}\quad 
        \overline{f} {\bm 1}_{[a,b]} = {\bm 1}_{[a,b]}, 
    \end{equation*}
    and such that ${\rm supp}(\underline{f_\eps})\subset [a+\eps,b-\eps]$ and ${\rm supp}(\overline{f_\eps})\subset [a-\eps,b+\eps]$.
    We also denote by $\Pi_{[a,b]}$ the spectral projector of $-\hbar^2\Delta_{sR}$ on the energy window $[a,b]$. Then,we can write
    \begin{equation*}
        \sum_{a\leq \delta_k(\hbar)\leq b}\langle B\varphi_k,\varphi_k\rangle = {\rm Tr}\left( B \,\Pi_{[a,b]}\right) = {\rm Tr}\left(\underline{f_\eps}(-\hbar^2\Delta_{sR})B\right)+ {\rm Tr}\left((\overline{f_\eps}(1-\underline{f_\eps}))(-\hbar^2\Delta_{sR}) B\,\Pi_{[a,b]}\right).
    \end{equation*}
    The second term can be estimated using the semiclassical functional calculus given by Theorem~\ref{thm:functionalcalc} and Lemma~\ref{lem:traceclass}, from which we can deduce that
    \begin{equation*}
        {\rm Tr}\left(\left|\overline{f_\eps}(1-\underline{f_\eps}))(-\hbar^2\Delta_{sR})\right|\right)\leq (c_1(\eps) + c_2(\hbar,\eps))\hbar^{-Q},
    \end{equation*}
    where $\lim_{\eps\rightarrow 0}c_1(\eps) = 0$ by support condition on the principal symbol of this operator relatively to $\eps$, following the definition of $\underline{f_\eps}$ and $\overline{f_\eps}$, and where $\lim_{\hbar\rightarrow 0}c_2(\hbar,\eps) = 0$.

    By the same argument applied to the first term, we readily get that
    \begin{equation*}
        \begin{aligned}
            \hbar^Q {\rm Tr}\left(B \,\Pi_{[a,b]}\right)
            &= \hbar^Q {\rm Tr}\left(\underline{f_\eps}(-\hbar^2\Delta_{sR})B\right) + \O\left(\|B\,\Pi_{[a,b]}\|_{\mathcal{L}(L^2(M))}(c_1(\eps)+c_2(\hbar,\eps)\right)\\
            &= \int_{\Ghat M}{\rm Tr}_{\mathcal{H}_\pi}\left(\underline{f_\eps}(H)(x,\pi)\beta(x,\pi)\right)d\hat{\mu}_{\Ghat M}(x,\pi) + \O_\eps(\hbar) + \O(c_1(\eps)+c_2(\hbar,\eps)).
        \end{aligned}
    \end{equation*}
    Thus, we have
    \begin{equation*}
        \limsup_{\hbar\rightarrow 0} \hbar^Q {\rm Tr}\left(B \,\Pi_{[a,b]}\right) = \int_{\Ghat M}{\rm Tr}_{\mathcal{H}_\pi}\left(\underline{f_\eps}(H)(x,\pi)\beta(x,\pi)\right)d\hat{\mu}_{\Ghat M}(x,\pi) + \O(c_1(\eps)),
    \end{equation*}
    and similarly for $\liminf_{\hbar\rightarrow 0}$. We conclude by dominated convergence letting $\eps$ go to 0.
\end{proof}

\end{document}
